% !PRINTED 118
\weberpage{118}
%<<CONT

\[
 y=\frac{-(u+v)\pm\sqrt{-3(u+v)^2+12uv}}{2}
  =\frac{-(u+v)\pm\sqrt{-3}(u-v)}{2}.
\]
Setzen wir also
\begin{equation}
\varepsilon=\frac{-1+\sqrt{-3}}{2},
\tag{11}
\end{equation}
woraus
\begin{equation}
\varepsilon^3=1,
\qquad \varepsilon^2+\varepsilon+1=0,
\qquad
\varepsilon^2=\frac{-1-\sqrt{-3}}{2}
\tag{12}
\end{equation}
folgt, so erhält man die drei Wurzeln der cubischen Gleichung (6)
% [sic] Source-original mathematical cross-reference: the printed text says (6), although the cubic equation itself is (1).
\begin{equation}
\begin{array}{rcl}
\alpha&=&u+v,\\
\beta&=&\varepsilon u+\varepsilon^2v,\\
\gamma&=&\varepsilon^2u+\varepsilon v.
\end{array}
\tag{13}
\end{equation}
Darin ist $\varepsilon$ eine von $1$ verschiedene dritte Einheitswurzel, die hier ohne die trigonometrischen Functionen algebraisch ausgedrückt ist. Setzt man $\varepsilon u$ oder $\varepsilon^2u$ für $u$, also $\varepsilon^2v$, $\varepsilon v$ für $v$, so vertauschen sich die drei Grössen $\alpha,\beta,\gamma$ unter einander.

Der Ausdruck (9) für die Wurzel einer cubischen Gleichung wird die \gesperrt{Cardanische Formel} genannt.

\begin{center}
{\large \S.\ 36.}

\vspace{0.5em}
{\large\gesperrt{Der Cayley'sche Ausdruck der Cardanischen Formel.}}
\end{center}

Jede Cubikwurzel hat, wie wir gesehen haben, drei verschiedene Werthe, die man aus einem von ihnen erhält durch Multiplication mit $1,\varepsilon,\varepsilon^2$. So hat also auch jede der beiden Grössen $u,v$, wie sie durch (8), \S.\,35 definirt sind, drei Werthe und die Summe $u+v$ hat also, wenn man nur die Bedingungen (8) berücksichtigt, \gesperrt{neun verschiedene} Werthe; von diesen geben aber nur drei Lösungen der cubischen Gleichung, und erst durch Zuziehung der Relation (4), \S.\,35
\begin{equation}
3uv=-a
\tag{1}
\end{equation}
werden unter den neun Werthen die brauchbaren ausgesondert.
